\section{Discussion and Limitations}
\label{sec:discussion}
\subsection{Collection Cost and Deployment Scope}
The synthetic replay reports 13,192.4 events/sec and 6.91 MB on a fixed stream. Those measurements do not establish zero interference with domain authentication, bounded history storage, or end-to-end CA polling cost. The lineage reference recomputes retained evidence, and the fault-budget reference enumerates deletion hypotheses. The recorded live command envelopes and offline inference timings quantify two parts of this cost: median recorded trial envelopes are 50.91~s and 32.14~s in the two environments, while median offline identity reasoning is 10.99~ms and 9.18~ms. A three-pair collection pilot measures periodic export/reference updates and verifies 48 legitimate authentications. The paired median client process-envelope increases are 48.6--100.3~ms, with sampled Python collector RSS of 56.5--59.5~MiB, excluding children and the VM. This quantifies a short isolated-lab workload; continuous high-load operation, lifetime retention and production interference remain unestablished.

\subsection{Hybrid Identity Evaluation}
The experiments cover on-premises ADCS and KDC records. Extending the evidence contract to Microsoft Entra ID would require additional identity-mapping and audit-source assumptions. Cross-tenant tokens and synchronization-account compromise are outside the evaluated workload. They should be treated as separate validation tasks before making claims about hybrid identity protection.

\subsection{Threats to Validity and Limitations}
We distinguish the following threats to validity and experimental limitations:
\begin{enumerate}
    \item \textbf{Evidence Scale and Label Validity:} The archive's scenario names do not establish seven successful attacks. Four final fixed-request matrices supply 42 successful exact-certificate authentications and six certificate-bound rejections, with nine privileged uses and matched self-use controls. A prospective transfer cohort adds 24 distinct certificates, including four further privileged uses: the combined primary studies contain 72 certificate cases, 66 successful uses and six rejections. The transfer cohort has new identity namespaces and concurrent multi-template activity but shares the Windows image family. They establish controlled outcomes, not enterprise prevalence. Independent Windows 2019 builds 17763.4974 and 17763.9376 are measured. Current-patch controls reveal request-form differences, but the additional post-rollback workflow prevents attributing every old/new difference solely to the update. Eighteen missing contemporaneous CA exports are explicitly supplemented by a later read-only snapshot; original captures remain intact. Original source images remain read-only.
    \item \textbf{Static Model and Generator Validity:} The matched audit checks reachable-path labels and checks 6,000 checkpoint predictions. Repeated training/validation archetypes, fixed first-step transforms, hand-initialized projections, and the distinction between feature-channel and bilinear gating restrict what its numbers establish. A separate literal bilinear study removes those implementation and repeated-whole-input problems with trainable encoders and typed family-disjoint splits. Both literal bilinear and additive attention still miss all 450 long-path positives per arm, without a classification advantage for gating. A verifier restricted to neural positives cannot repair these misses, and symbolic label agreement is not independent Windows ground truth.
    \item \textbf{Parametric Disruption Costs vs. Interactive Administrator Feedback:} In the budgeted interdiction model, operational disruption costs $c(e)$ are assigned parametrically according to relation type ($c(\text{MemberOf}) \in [2.8, 3.8], c(\text{GenericAll}) \in [2.0, 3.0], c(\text{Enroll}) \in [1.0, 1.4]$) and operational edge immutability constraints ($E_b \subsetneq E$). In production enterprise deployments, operational criticality is context-sensitive and evolving. The exact small-instance audit gives concrete optimum comparisons rather than a universal heuristic guarantee. A metadata-preserving rerun also checks all 120 original graph/budget instances; it changes nine objective outcomes and 24 cut counts. The primary table and revoked-edge comparison use full edge-attribute restoration (Appendix~\ref{sec:artifact_corrections}). Eliciting and validating administrator-assigned costs requires a separate evaluation; the graph-hardening study~\cite{goel2024hardening} is not cited as evidence of such preference elicitation.
    \item \textbf{Trust, Time, and Necessary Causation:} Candidate retention under a bridge fault budget assumes correct raw issuance/KDC records and captured true issuance. Consistently forged raw histories defeat the methods, while faults exceeding the declared budget invalidate its retention guarantee. A counterexample shows that self-consistent false order can also produce a wrong definite overlap. A grouped-source reference handles the six-record correlated-export control under an explicit one-source assumption, which differs from a record budget. The operation-order pilot uses the controlled orchestrator as an additional trusted source and does not repair the failed Windows audit-clock model. Temporal overlap alone does not establish necessity. The matched CSR/key interventions support a conditional contrast under the recorded policy fields and controlled CA refresh, while unobserved concurrent changes and missing restoration audit events remain distinct evidence concerns.
\end{enumerate}

\section{Conclusion}
\label{sec:conclusion}
CertGraph links certificate issuance and use under an explicit issuer-scoped evidence contract. Its live evaluation verifies 66 successful certificate authentications and six certificate-bound rejections, including 13 administrative uses. The current-patch controls show why a restored template scan and a certificate-use investigation answer different questions. The identity contract can also be implemented using guarded SQL; the evaluated contribution concerns its security requirements and evidence handling rather than a unique reasoning primitive.

The graph-learning results do not establish a classification advantage for bilinear gating. Both matched literal models miss the tested long-path positives, so a verifier that examines only neural-positive cases remains incomplete for this workload. The remediation experiments also show heuristic optimality gaps and budget-dependent differences from greedy selection. Larger deployments require separate measurements of collection cost, retained history and administrator-assigned disruption costs. The present results support these specific conclusions within the recorded laboratory and generator conditions.

\section*{Author Contributions}
\noindent\textbf{Md.~Hasan~Shahriar}: Conceptualization, Methodology, Software, Formal Analysis, Investigation, Writing -- Original Draft, Visualization, Project Administration.
\textbf{Taufiq~Biswas}: Validation, Data Curation, Formal Analysis, Writing -- Review \& Editing.
\textbf{Md.~Ramim~Hossain}: Validation, Resources, Investigation, Writing -- Review \& Editing.

\section*{Conflicts of Interest}
\noindent The authors declare that they have no known competing financial interests or personal relationships that could have appeared to influence the work reported in this paper.

\section*{Funding}
\noindent This research received no specific grant from any funding agency.

\section*{Ethics and Dual-Use Considerations}
\noindent The reported interventions concern isolated, researcher-controlled Windows forests and laboratory accounts. They test certificate abuse and permission changes solely to assess defensive evidence and remediation. Original disk images were retained read-only, and writable overlays were used for the experiments. No production organization or unsuspecting user was part of the evaluated workload. Credential-bearing material, private keys and reusable authentication tokens are excluded from the review artifact. The review package retains synthetic controls, verification code and derived results; the original live evidence remains in controlled storage.

\section*{Data and Code Availability}
\noindent The supplementary review artifact contains the synthetic study inputs, source snapshots, verification results and manuscript evidence tables described in this paper. The artifact README identifies omitted credential-bearing laboratory material and the resulting reproduction limits. The project repository is \url{https://github.com/hasanshahriar32/CertGraph-CIPG}; a persistent, versioned public archive of this revision has not yet been deposited. No archive DOI is claimed.
